% Part: incompleteness
% Chapter: incompleteness-provability
% Section: first-incompleteness-thm

\documentclass[../../../include/open-logic-section]{subfiles}

\begin{document}

\olfileid{inc}{inp}{1in}

\olsection{پہلا نامکملیت قضیہ}

ہن اسیں پہلے نامکملیت قضیے دا گوڈل والا اصل ثبوت بیان
کر سکدے آں۔ فرض کرو $\Th{T}$ ایہو جیہا نظریہ اے جیہدی
مسلّمہ بندی قابلِ حساب اے، تے جیہدی زبان حساب دی زبان دی
توسیع اے۔ ایہ وی فرض کرو کہ $\Th{T}$ وچ $\Th{Q}$ دے
مسلّمے شامل نیں۔ خاص طور اُتے ایس توں نکلدا اے کہ
$\Th{T}$ قابلِ حساب فنکشناں تے تعلقاں دی نمائندگی کردا اے۔

اسیں ویکھیا اے کہ جے فارمولیاں تے ثبوتاں نوں عدداں نال
کوڈ کرن دا معقول طریقہ ہووے، تاں تعلق $\Prf[\Th{T}](x,y)$
قابلِ حساب اے۔ $\Prf[\Th{T}](x,y)$ اُدوں تے صرف اُدوں
قائم اے جدوں $x$ اوس !!a{derivation} دا گوڈل نمبر ہووے
جیہڑا گوڈل نمبر~$y$ والے !!{formula} دا~$\Th{T}$ وچ
اخذ اے۔ گوڈل نے اپنے خاص نظریے لئی اپنے مقالے وچ دتی 45
فنکشناں تے تعلقاں دی فہرست ورت کے وکھایا کہ ایہ تعلق بنیادی
بازگشتی اے۔ اوس فہرست دا 45واں تعلق، $x B y$، اوہی
$\Prf[\Th{T}](x,y)$ اے جیہڑا گوڈل دے خاص نظریے
~$\Th{T}$ لئی ورتیا جاندا اے۔ چیتے رکھو،
گوڈل اپنے مقالے وچ جیہنوں “بازگشتی” آکھدا اے، اسیں ہن
اوہنوں “بنیادی بازگشتی” آکھدے آں۔

کیوں جے $\Prf[\Th{T}](x,y)$ قابلِ حساب اے، ایس لئی
$\Th{T}$ وچ اوہدی نمائندگی ہو سکدی اے۔ اسیں اوہدی
نمائندگی کرن والے فارمولے نوں $\OPrf[\Th{T}](x,y)$ آکھاں
گے۔ $\OProv[\Th{T}](y)$ نوں فارمولا
$\lexists[x][\OPrf[\Th{T}](x,y)]$ منّ لئیے۔ ایہ گوڈل دی
فہرست دے 46ویں تعلق، $\fn{Bew}(y)$، نوں بیان کردا اے۔
گوڈل آکھدا اے کہ ایہ اکّو تعلق اے جیہدے بارے پکے طور اُتے
نہیں آکھیا جا سکدا کہ ایہ بازگشتی اے۔ اوہدا مطلب غالباً ایہ سی
کہ تعریف توں ایہ صاف نہیں کہ تعلق قابلِ حساب اے یا نہیں۔
پچھوں ہون والی پیش رفت توں پتا لگدا اے کہ ایہ قابلِ حساب نہیں۔

فرض کرو $\Th{T}$ اک !!{axiomatizable} نظریہ اے جیہدے وچ
~$\Th{Q}$ شامل اے۔ تاں $\Prf[\Th{T}](x, y)$ قابلِ فیصلہ
اے، ایس لئی~$\Th{Q}$ وچ !!a{formula}~$\OPrf[\Th{T}](x, y)$
نال اوہدی نمائندگی ہو سکدی اے۔ $\OProv[\Th{T}](y)$ اوہی
فارمولا اے جیہڑا اسیں اُتے بیان کیتا۔ مقرر نقطے دے لمے نال
اک فارمولا $!G_\Th{T}$ ملدا اے جیہدے لئی $\Th{Q}$،
تے ایس لئی $\Th{T}$ وی، ایہ !!{derive}s کردا اے،
\begin{equation}
\ollabel{eqn:qpf}
!G_\Th{T} \liff \lnot \OProv[\Th{T}](\gn{!G_\Th{T}}).
\end{equation}
دھیان دیو کہ $!G_\Th{T}$ دا مطلب ایہ اے، “$!G_\Th{T}$،
~$\Th{T}$ وچ !!{derivable} نہیں۔”

\begin{lem}\ollabel{lem:cons-G-unprov}
جے $\Th{T}$،~$\Th{Q}$ دی توسیع والا اک بے تضاد
!!{axiomatizable} نظریہ اے، تاں $\Th{T} \Proves/ !G_\Th{T}$۔
\end{lem}

\begin{proof}
فرض کرو $\Th{T}$، $!G_\Th{T}$ نوں !!{derive}s کردا اے۔
تاں اک !!a{derivation} \emph{ضرور موجود اے}۔ ایس لئی کسے عدد
$m$ لئی تعلق $\Prf[\Th{T}](m,
\Gn{!G_\Th{T}})$ قائم اے۔ فیر $\Th{Q}$، جملہ
$\OPrf[\Th{T}](\num m, \gn{!G_\Th{T}})$ !!{derive}s کردا اے۔
سو $\Th{Q}$ ایہ وی !!{derive}s کردا اے،
$\lexists[x][\OPrf[\Th{T}](x,\gn{!G_\Th{T}})]$، جیہڑا تعریف
مطابق $\OProv[\Th{T}](\gn{!G_\Th{T}})$ اے۔ \olref{eqn:qpf}
نال $\Th{Q}$، $\lnot !G_\Th{T}$ !!{derive}s کردا اے۔
کیوں جے $\Th{T}$ وچ $\Th{Q}$ شامل اے، ایس لئی
~$\Th{T}$ وی ایہ اخذ کردا اے۔ اسیں وکھایا اے کہ جے
$\Th{T}$، $!G_\Th{T}$ !!{derive}s کرے، تاں اوہ
$\lnot !G_\Th{T}$ وی !!{derive}s کرے گا۔ ایس طرح اوہ متضاد
ہو جائے گا۔
\end{proof}

\begin{defn}
\ollabel{thm:oconsis-q}
نظریہ $\Th{T}$، \emph{$\omega$-بے تضاد} اے جے ایہ شرط
پوری ہووے۔ جے $\lexists[x][!A(x)]$ کوئی وی جملہ اے تے
$\Th{T}$، $\lnot
!A(\num 0)$، $\lnot !A(\num 1)$، $\lnot !A(\num 2)$،
\dots سبھ !!{derive}s کردا اے، تاں $\Th{T}$،
$\lexists[x][!A(x)]$ ثابت نہیں کردا۔
\end{defn}

دھیان دیو کہ ہر $\omega$-بے تضاد نظریہ عام بے تضاد وی اے۔
وجہ ایہ اے کہ جے $\Th{T}$ متضاد اے، تاں $\Th{T} \Proves !A$
ہر~$!A$ لئی۔ خاص طور اُتے، جے $\Th{T}$ متضاد اے، تاں
اوہ $\lnot !A(\num n)$ ہر~$n$ لئی !!{derive}s کردا اے
تے~$\lexists[x][!A(x)]$ وی !!{derive}s کردا اے۔ سو جے
$\Th{T}$ متضاد اے، تاں اوہ $\omega$-بے تضاد نہیں۔ ایس
شرط دے عکسِ نقیض توں نکلدا اے کہ جے $\Th{T}$،
$\omega$-بے تضاد اے، تاں اوہ لازماً عام بے تضاد اے۔

\begin{lem}\ollabel{lem:omega-cons-G-unref}
جے $\Th{T}$ اک $\omega$-بے تضاد !!{axiomatizable} نظریہ اے
جیہڑا~$\Th{Q}$ دی توسیع اے، تاں
$\Th{T} \Proves/ \lnot !G_\Th{T}$۔
\end{lem}

\begin{proof}
اسیں وکھاواں گے کہ جے $\Th{T}$، $\lnot !G_\Th{T}$
!!{derive}s کرے، تاں اوہ $\omega$-بے تضاد نہیں۔ فرض کرو
$\Th{T}$، $\lnot !G_\Th{T}$ !!{derive}s کردا اے۔ جے
$\Th{T}$ متضاد اے، تاں اوہ $\omega$-بے تضاد نہیں تے
دلیل پوری اے۔ دوجی صورت وچ $\Th{T}$ بے تضاد اے، ایس لئی
اوہ $!G_\Th{T}$ !!{derive} نہیں کردا، \olref{lem:cons-G-unprov}
مطابق۔ کیوں جے $!G_\Th{T}$ دا کوئی !!{derivation}
$\Th{T}$ وچ موجود نہیں، ایس لئی $\Th{Q}$ ایہ سبھ
!!{derive}s کردا اے،
\[
\lnot \OPrf[\Th{T}](\num 0, \gn{!G_\Th{T}}), \lnot \OPrf[\Th{T}](\num 1,
\gn{!G_\Th{T}}), \lnot \OPrf[\Th{T}](\num 2, \gn{!G_\Th{T}}), \dots
\]
تے~$\Th{T}$ وی ایہ سبھ اخذ کردا اے۔ دوجے پاسے
\olref{eqn:qpf} مطابق $\lnot
!G_\Th{T}$، $\lexists[x][\OPrf[\Th{T}](x,\gn{!G_\Th{T}})]$
دے ہم ارز اے۔ ایس لئی $\Th{T}$، $\omega$-بے تضاد نہیں۔
\end{proof}

\begin{prob}
ہر $\omega$-بے تضاد نظریہ عام بے تضاد اے۔ وکھاؤ کہ الٹی
شرط لازم نہیں، یعنی کجھ نظریے بے تضاد نیں پر
$\omega$-بے تضاد نہیں۔ ایس لئی وکھاؤ کہ $\Th{Q} \cup
\{\lnot !G_\Th{Q}\}$ عام بے تضاد اے پر $\omega$-بے تضاد
نہیں۔
\end{prob}

\begin{thm}
\ollabel{thm:first-incompleteness}
فرض کرو $\Th{T}$ کوئی وی $\omega$-بے تضاد
!!{axiomatizable} نظریہ اے جیہڑا~$\Th{Q}$ دی توسیع اے۔
تاں $\Th{T}$ مکمل نہیں۔
\end{thm}

\begin{proof}
جے $\Th{T}$، $\omega$-بے تضاد اے، تاں اوہ عام بے تضاد
وی اے۔ ایس لئی $\Th{T}
\Proves/ !G_\Th{T}$، \olref{lem:cons-G-unprov} مطابق۔
\olref{lem:omega-cons-G-unref} مطابق
$\Th{T} \Proves/ \lnot !G_\Th{T}$۔ ایس توں نکلدا اے کہ
$\Th{T}$ نامکمل اے، کیوں جے اوہ نہ $!G_\Th{T}$ !!{derive}s
کردا اے تے نہ $\lnot !G_\Th{T}$۔
\end{proof}

\end{document}
